\begin{center}\begin{tikzpicture}[font=\small,>=stealth,line join=round]
\draw[gray!40,rounded corners=3pt] (0,0) rectangle (14.2,8.5);\node at (7.1,7.9) {代入零以后，函数值必须是自己的相反数};
\draw[->,gray] (1,5.5)--(13.1,5.5);\foreach \x/\v in {3/-2,7.1/0,11.2/2}{\draw[gray] (\x,5.5)--(\x,5.35);\node[below=6pt] at (\x,5.5){$\v$};}\fill[blue] (3,5.5) circle(2.3pt);\fill[orange!85!black] (11.2,5.5) circle(2.3pt);\draw[->,blue] (11.2,5.7) .. controls (10,6.8) and (4,6.8) .. (3,5.7);\node at (7.1,6.9){取相反数：$2\mapsto-2$};\node at (7.1,4.35){$2$ 与 $-2$ 是两个不同的数};\draw[->,gray] (1,2.5)--(13.1,2.5);\fill[orange!85!black] (7.1,2.5) circle(2.5pt);\node[below=5pt] at (7.1,2.5){$0$};\draw[->,blue] (7.4,2.6) .. controls (9,3.7) and (5.2,3.7) .. (6.8,2.6);\node at (7.1,1.4){只有 $0$ 取相反数后仍在原位置：$0=-0$};
\node[text width=12.8cm,align=center] at (7.1,.65) {两条数轴比例相同；$f(0)=-f(0)$ 正是下方零点的自反关系。};
\end{tikzpicture}\end{center}
